\documentclass{article}
\usepackage{tikz}
\usepackage{pgfplots}
\usepackage{amsmath, amsfonts}
\usepackage{enumerate}
\usepackage{makecell}
\usepackage{xypic}
\usepackage{fancyvrb}
\usepackage{listings}
\input{title.tex}

\newcommand{\N}{\mathbb{N}}
\newcommand{\R}{\mathbb{R}}
\begin{document}
    \maketitle
    \section{}
    \section{}
    \begin{enumerate}[a)]
        \item
            Wir zeigen CSP $\leq_p$ Rucksack.
            Wir nehmen einen Rucksack mit gegenständen $I=\{1,\hdots,n\}$,
            einer Wertefunktion $w:I\rightarrow \N$ und einer
            Gewichtefunktion $g:I\rightarrow \N$.
            Ebenso gibt es ein Maximalgewicht $\mathbb G\in\N$.
            Anhand dieses Rucksacks konstruieren wir

            \begin{itemize}
                \item einen Graph $G(V, E)$
                    mit $V=\{v_1,\hdots v_n\}$
                    und $E=\{v_i,v_j\}_{i>j\in I}$,
                    Damit wird sichergestellt, dass jeder Knoten nur
                    einmal besucht werden kann.
                    $v_1$ ist Start- und $v_n$ Endknoten.
                \item eine Recoursefunktion $r:E\rightarrow R^+$
                    mit $r(v_i)=g(i)$ für $v_i\in V$
                \item eine Kostenfunktion
                    $c(v_i)=-w(i)$ für $v_i\in V$
                \item Eine Maximalresource $R=\mathbb G$
            \end{itemize}
            Wenn wir nun Annahmen, dass CSP nicht NP-Schwer ist, dann Lösen wir
            CSP für den in Linearer zeit Konstruierten Graph und erhalten
            einen Pfad $\pi$ mit
            $\sum_{e\in\pi} c(e)$ minimal und $\sum_{e\in\pi}r(e)\leq R$.
            Somit haben wir Elemente $I'=\{i|v_i\in e\}\subseteq I$
            mit
            $\sum_{i\in I'} w(i)$ maximal und $\sum_{i\in I'}g(i)\leq G$
            die das Rucksackproblem in Polynomischer lösen. Da Rucksack
            aber NP-hart ist, liegt hier ein Widerspruch zur Annahme vor.
        \newpage
        \item
        \begin{lstlisting}[mathescape=true]
func CSP( graph G(V, E) , source , destination , R) {
    // this holds the shortest from the source to
    // node v with constraint $\leq R$ with all initially at $\infty$
    for $v\in V, r\leq R$ {
        min_distange[$v$][$r$] := $\infty$;
        min_distange[source][0] := 0;
    }
    // iterating over all constraints
    for $k\in \{0,\hdots,R\}$ {
        for $u \in V$ {
        // checking if $u$ is reachable
            if ( min_distange[$u$][$k$] < $\infty$) {
                // checking all outgoing vertices
                for each outgoing edge v {
                    weight := $w(u , v)$;
                    // reassign if lower
                    if $(k + weight \leq R)$ {
                    min_distange[v][k+weight] = min(
                        min_distange[v][k+weight],
                        min_distange[u][k] + weight
                    );
                }
            }
        }
    }
    // checking shortest path for all constraints
    result := $\infty$
    for $k\in\{0,\hdots, R\}$ {
        result = min(result, min_distange[target][k]);
    }
    if (result = $\infty$) {
        return NO PATH FOUND;
    }
    else {
        return result;
    }
}
        \end{lstlisting}
    \end{enumerate}
\end{document}
